% SEGMENT OLP-0444-S01
% Part: normal-modal-logic
% Chapter: completeness
% Section: lindenbaums-lemma

\documentclass[../../../include/open-logic-section]{subfiles}

\begin{document}

\olfileid{nml}{com}{lin}

\olsection{லிண்டன்பாமின் துணைத்தேற்றம்}

ஒவ்வொரு $\Sigma$-முரணற்ற !!{formula}s கணமும் குறைந்தது ஒரு
\emph{நிறைவான} $\Sigma$-முரணற்ற கணத்தில் அடங்கியுள்ளது என்பதை
லிண்டன்பாமின் துணைத்தேற்றம் நிறுவுகிறது. நிறைவான ஒவ்வொரு
$\Sigma$-முரணற்ற கணம்~$\Delta$ க்கும்,~$\Delta$ விலுள்ள !!{formula}s
அனைத்தும், அவை மட்டுமே, மெய்யாகும் ஓர் உலகம் நியம மாதிரியில் உள்ளது என்பதை
நமது நியம மாதிரிக் கட்டமைப்பு காட்டும். ஆகவே ஒவ்வொரு
$\Sigma$-முரணற்ற கணமும் நியம மாதிரியில் ஏதோவோர் உலகில் மெய்யாகும் என்பதற்கு
லிண்டன்பாமின் துணைத்தேற்றம் உத்தரவாதமளிக்கிறது.

\begin{thm}[லிண்டன்பாமின் துணைத்தேற்றம்]\ollabel{thm:lindenbaum}
  $\Gamma$ ஆனது~$\Sigma$-முரணற்றது எனில்,~$\Gamma$ வை விரிவாக்கும்
  நிறைவான $\Sigma$-முரணற்ற கணம்~$\Delta$ ஒன்று உள்ளது.
\end{thm}

\begin{proof}
மொழியின் எல்லா வாய்பாடுகளையும் $!A_0$, $!A_1$, \dots{} என முழுமையாகப்
பட்டியலிடுக (மீள்கூறல்கள் அனுமதிக்கப்படுகின்றன). எடுத்துக்காட்டாக, முதலில்
$\Obj p_0$ ஐப் பட்டியலிட்டு, ஒவ்வொரு படிநிலை~$n \ge 1$ இலும்
$\Obj p_0$, \dots,~$\Obj p_n$ ஆகிய மாறிகளை மட்டும் பயன்படுத்தும்
நீளம்~$n$ உள்ள இறுதிக்குட்பட்ட எண்ணிக்கையிலான வாய்பாடுகளைப் பட்டியலிடுக.
!!{formula}s கொண்ட கணங்கள்~$\Delta_n$ ஐ~$n$ மீதான தொகுத்தறிதலால்
வரையறுத்து, பின்னர் $\Delta = \bigcup_n \Delta_n$ என வைக்கிறோம்.
முதலில் $\Delta_0 = \Gamma$ என வைக்கிறோம். $\Delta_n$
வரையறுக்கப்பட்டுள்ளது எனக் கொண்டு, $\Delta_{n+1}$ ஐப் பின்வருமாறு
வரையறுக்கிறோம்:
\[
\Delta_{n+1} =
\begin{cases}
  \Delta_n \cup \{!A_n\}, & \text{$\Delta_n \cup \{ !A_n\}$
    ஆனது~$\Sigma$-முரணற்றது எனில்;} \\
  \Delta_n \cup \{ \lnot !A_n\}, & \text{இல்லையெனில்.}
\end{cases}
\]
இப்போது $\Delta = \bigcup_{n=0}^\infty \Delta_n$ என்க.

இவ்வரையறை வேண்டிய பண்புகளைக் கொண்ட ஒரு கணம்~$\Delta$ ஐ உண்மையில்
தருகிறது, அதாவது $\Gamma \subseteq \Delta$ என்றும்~$\Delta$ நிறைவான
$\Sigma$-முரணற்ற கணம் என்றும் காட்ட வேண்டும்.

கட்டமைப்பின்படி $\Delta_0 \subseteq \Delta$, மேலும் $\Delta_0 = \Gamma$
என்பதால், $\Gamma \subseteq \Delta$ என்பது தெளிவு. உண்மையில், ஒவ்வொரு
$n$ க்கும் $\Delta_n \subseteq \Delta$; ஏனெனில்~$\Delta$ ஆனது எல்லா
$\Delta_n$ களின் ஒன்றிப்பு. (கட்டமைப்பின் ஒவ்வொரு படியிலும், ஏற்கெனவே
கட்டமைக்கப்பட்ட கணத்திற்கு !!a{formula} ஒன்றைச் சேர்ப்பதால்,
$\Delta_n \subseteq \Delta_{n+1}$; ஆகவே~$\subseteq$ கடப்புப் பண்புடையது
என்பதால், $n \le m$ ஆகும்போதெல்லாம் $\Delta_n \subseteq \Delta_{m}$.)
கட்டமைப்பின் ஒவ்வொரு படிநிலையிலும் $!A_n$ அல்லது $\lnot !A_n$ ஐச்
சேர்க்கிறோம்; மேலும் ஒவ்வொரு !!a{formula} உம் எல்லா~$!A_n$ களின்
பட்டியலில் (குறைந்தது ஒருமுறை) வருகிறது. ஆகவே, ஒவ்வொரு~$!A$ க்கும்
$!A \in \Delta$ அல்லது $\lnot !A \in \Delta$; எனவே வரையறைப்படி
$\Delta$ நிறைவானது.

இறுதியாக, $\Delta$ ஆனது~$\Sigma$-முரணற்றது என்று காட்ட வேண்டும்.
இதற்காக, (a)~$\Delta$ ஆனது~$\Sigma$-முரணுடையதாக இருந்தால் ஏதோவொரு
$\Delta_n$ உம்~$\Sigma$-முரணுடையதாக இருக்கும் என்றும்,
(b)~எல்லா~$\Delta_n$ களும்~$\Sigma$-முரணற்றவை என்றும் காட்டுகிறோம்.

$\Delta$ ஆனது~$\Sigma$-முரணுடையதாக இருந்தது என்க. அப்போது
$\Delta \Proves[\Sigma] \lfalse$; அதாவது, $!A_1$, \dots,~$!A_k \in
\Delta$ ஆகவும், $\Sigma \Proves !A_1 \lif (!A_2 \lif \cdots (!A_k
\lif \lfalse)\dots)$ ஆகவும் இருக்கும் வாய்பாடுகள் உள்ளன.
$\Delta = \bigcup_{n=0}^\infty \Delta_n$ என்பதால், ஒவ்வொரு
$!A_i \in \Delta_{n_i}$ ஆகவும் இருக்கும் ஏதோவோர்~$n_i$ உள்ளது.
இவற்றில் மிகப்பெரியதை~$n$ என்க. $n_i \le n$ என்பதால்,
$\Delta_{n_i} \subseteq \Delta_n$. எனவே எல்லா~$!A_i$ களும் ஏதோவொரு
$\Delta_n$ இல் உள்ளன. இது $\Delta_n \Proves[\Sigma] \lfalse$,
அதாவது $\Delta_n$ ஆனது~$\Sigma$-முரணுடையது, என்று பொருள்படும்.

ஒவ்வொரு $\Delta_n$ உம்~$\Sigma$-முரணற்றது என்று காட்ட,~$n$ மீது ஓர்
எளிய தொகுத்தறிதலைப் பயன்படுத்துகிறோம். $\Delta_0 = \Gamma$; மேலும்
$\Gamma$ ஆனது~$\Sigma$-முரணற்றது என்று கருதினோம். ஆகவே கூற்று~$n = 0$
க்கு உண்மை. இப்போது அது~$n$ க்கு உண்மை, அதாவது $\Delta_n$ ஆனது
$\Sigma$-முரணற்றது, என்க. $\Delta_n \cup \{!A_n\}$ ஆனது
$\Sigma$-முரணற்றது எனில், $\Delta_{n+1}$ அந்தக் கணமே; இல்லையெனில் அது
$\Delta_n \cup \{\lnot!A_n\}$. முதல் நேர்வில், $\Delta_{n+1}$
தெளிவாகவே $\Sigma$-முரணற்றது. எனினும்,
\olref[prf][con]{prop:consistencyfacts}\olref[prf][con]{prop:consistencyfacts-c}
படி, $\Delta_n \cup \{!A_n\}$ அல்லது $\Delta_n \cup \{\lnot!A_n\}$
ஆகியவற்றில் ஒன்று முரணற்றது; ஆகவே மற்ற நேர்விலும் $\Delta_{n+1}$
முரணற்றது.
\end{proof}

\begin{cor}\ollabel{cor:provability-characterization}
  $\Gamma \Proves[\Sigma] !A$ என்பது, அப்போதும் அப்போது மட்டுமே,
  $\Gamma$ வை விரிவாக்கும் ஒவ்வொரு நிறைவான $\Sigma$-முரணற்ற
  கணம்~$\Delta$ வுக்கும் $!A \in \Delta$ (இது $\Gamma = \emptyset$
  ஆகும்போதும் பொருந்தும்; அப்போது வாய்ப்புநிலை அமைப்பு~$\Sigma$ க்கான
  மற்றொரு சிறப்பிப்பு கிடைக்கிறது.)
\end{cor}

\begin{proof}
  $\Gamma \Proves[\Sigma] !A$ என்க; மேலும் $\Gamma$ வை விரிவாக்கும்
  எந்த நிறைவான $\Sigma$-முரணற்ற கணம்~$\Delta$ ஐயும் எடுத்துக்கொள்க.
  $!A \notin \Delta$ எனில், மீப்பெருமையால் $\lnot!A \in \Delta$;
  மேலும் $\Delta \Proves[\Sigma] !A$ (ஒருபோக்குத்தன்மையால்),
  $\Delta \Proves[\Sigma] \lnot!A$ (தற்சுட்டுத்தன்மையால்) ஆகியவை
  கிடைப்பதால், $\Delta$ முரணுடையது. மறுதிசையில், $\Gamma
  \Proves/[\Sigma] !A$ எனில், $\Gamma \cup \{ \lnot!A\}$ ஆனது
  $\Sigma$-முரணற்றது; லிண்டன்பாமின் துணைத்தேற்றத்தின்படி, $\Gamma \cup
  \{ \lnot!A \}$ ஐ விரிவாக்கும் நிறைவான முரணற்ற கணம்~$\Delta$ ஒன்று
  உள்ளது. முரணின்மையால், $!A \notin \Delta$.
\end{proof}

\end{document}
